\SourceLeaf{206}
durch Elemente in der zweiten Ebne dargestellt werden, die unter sich und, wie wir hinzufügen können, \textit{in demselben Sinn}, dieselben Winkel bilden, wie jene.

Wählt man für \(f\) eine linearische Function, so dass \(f\upsilon=A+B\upsilon\), wo die constanten Coëfficienten von der Form sind
\[
A=a+bi,\quad B=c+ei
\]
\noindent so wird
\[
\varphi\upsilon=B=c+ei
\]
\noindent also
\[
\sigma=\BareRoot{(cc+ee)},\quad\gamma=\mathrm{Arc.}\,\mathrm{tang}\frac{e}{c}
\]
\noindent Das Vergrösserungsverhältniss ist folglich in allen Punkten constant, und die Darstellung dem Dargestellten durchaus ähnlich.

Für jede andere Function \(f\) wird (wie man leicht beweisen kann) das Vergrösserungsverhältniss nicht constant sein, und die Aehnlichkeit also nur in den kleinsten Theilen Statt finden können.

Sind die Plätze, welche einer bestimmten Anzahl von gegebenen Punkten der ersten Ebne in der Darstellung entsprechen sollen, vorgeschrieben, so kann man leicht nach der gemeinen Interpolationsmethode die einfachste algebraische Function \(f\) finden, wodurch diese Bedingung erfüllt wird. Bezeichnet man nemlich die Werthe von \(x+iy\) für die gegebenen Punkte durch \(a,b,c\) u. s. w., und die correspondirenden Werthe von \(X+iY\) durch \(A,B,C\) u. s. w., so wird man
\[
\begin{aligned}
f\upsilon={}&\frac{(\upsilon-b)(\upsilon-c).\ .\ .}{(a-b)(a-c).\ .\ .}\mathord{.}A
+\frac{(\upsilon-a)(\upsilon-c).\ .\ .}{(b-a)(b-c).\ .\ .}\mathord{.}B\\
&+\frac{(\upsilon-a)(\upsilon-b).\ .\ .}{(c-a)(c-b).\ .\ .}\mathord{.}C+\text{etc.}
\end{aligned}
\]
\noindent setzen müssen, welches eine algebraische Function von \(\upsilon\) ist, deren Ordnung um eine Einheit kleiner ist, als die Anzahl der vorgegebenen Punkte. Für zwei Punkte, wo die Function linearisch wird, findet folglich vollkommene Aehnlichkeit Statt.

Man kann von diesem Verfahren in der Geodäsie eine nützliche Anwendung machen, um eine auf mittelmässige Messungen gegründete Karte, die im kleinen Detail gut, aber im Ganzen etwas verzerrt ist, in eine bessere zu verwandeln, wenn man die richtige Lage einer Anzahl von Punkten kennt. Es versteht sich jedoch, dass man bei einer solchen Umformung nicht viel über die Gegend hinausgehen darf, welche letztere Punkte umfassen.

Wenn man die \textit{zweite} Auflösung auf dieselbe Art durchführt, so findet man, dass der ganze Unterschied nur darin besteht, dass die Aehnlichkeit eine verkehrte ist, indem alle Elemente in der Darstellung zwar eben so grosse Winkel mit einander machen, wie im Dargestellten, aber in verkehrtem Sinn, so dass
